% Tableau des labs du cours INA dans NDAP SUKULU (propre au cours).
Le cours compte douze laboratoires, disponibles sous le code \codeUE. Les
onze premiers durent une heure et suivent l'ordre des chapitres~; le
douzième, le plus long et le plus détaillé, résume tout le cours~: il
couvre les études de cas intégratrices (chapitre~21), le benchmarking
(chapitre~25) et les fondements cryptographiques de la preuve (Partie~VI).

\begin{tableaucours}{lp{6.4cm}p{1.6cm}p{2.2cm}}{%
    \tableauHeader{Lab} &
    \tableauHeader{Sujet} &
    \tableauHeader{Chapitre} &
    \tableauHeader{Opération}
}
INA-LAB-01 & Forensique système~: timeline Linux et Windows & 14 & MVOM \\
INA-LAB-02 & Forensique réseau opérationnelle~: beaconing et courriel frauduleux & 15 & BAOBAB, NKAM \\
INA-LAB-03 & Journaux réseau sans capture~: DNS, proxy, pare-feu & 15 & MVOM \\
INA-LAB-04 & Forensique mobile~: SMS, appels, EXIF & 16 & BAOBAB, MBONGO \\
INA-LAB-05 & Forensique cloud~: CloudTrail et fenêtre aveugle & 17 & MVOM \\
INA-LAB-06 & Forensique IoT~: firmware de routeur et NVR & 18 & NYONG \\
INA-LAB-07 & OSINT~: domaines, profils, cryptoactifs & 19 & BAOBAB \\
INA-LAB-08 & Le rapport forensique~: relecture avant dépôt & 20 & MVOM \\
INA-LAB-09 & Anti-forensique~: timestomping et journaux effacés & 22 & MVOM \\
INA-LAB-10 & L'IA en forensique~: mesurer un outil, examiner des clips & 23 & MUNGO \\
INA-LAB-11 & Conteneurs~: audit Kubernetes et conteneur évincé & 24 & MVOM \\
INA-LAB-12 & Synthèse~: du dossier intégrateur à la preuve vérifiable (lab long) & 21, 25, 26--32 & Toutes \\
\end{tableaucours}
\vspace{-0.8em}
\begin{center}{\small\itshape Les laboratoires du cours dans NDAP SUKULU}\end{center}
